\paragraph{From short-horizon forecasting to route generation.}
Urban mobility traces record how people realize intent through space: not only \emph{where} they go, but \emph{how} they get there.
Recent progress in sequence models and diffusion models \citep{vaswani2017attention,ho2020ddpm} has made short-horizon trajectory prediction increasingly accurate.
Yet many practical and scientific tasks require a different capability: \emph{route generation} over a full trip.
Given an origin--destination intent and context (e.g., departure time and neighborhood semantics), a route generator must sample an entire path that is realistic, feasible, and diverse.

\paragraph{Why long-horizon route generation fails.}
We observe two recurrent failure modes when extending generative modeling from seconds to full trips.
First, intent can drift: purely autoregressive rollouts accumulate small errors and become globally inconsistent with the destination.
Second, feasibility can break: when emphasizing diversity (e.g., via latent-variable models \citep{kingma2014auto} or standard diffusion), generators often hallucinate physically invalid paths that cut through obstacles or deviate systematically from the road network---especially when explicit map supervision is absent or unreliable.

\paragraph{Root cause: topological multi-modality and destructive averaging.}
The core difficulty is \emph{topological multi-modality}.
Even for the same origin and destination, multiple corridor-level choices can be equally reasonable.
End-to-end coordinate generators trained with likelihood or Euclidean losses therefore face a destructive ``averaging'' pressure: the mean of several feasible routes is often an infeasible straight line.
This suggests that route generation should be modeled as a \emph{decision--execution} process rather than a single homogeneous sequence problem.

\paragraph{A diagnostic baseline: averaging vs.\ blur in continuous output space.}
We first provide a diagnostic experiment (Fig.~\ref{fig:mode-collapse}) using end-to-end baselines trained under the same data and conditioning.
To avoid cherry-picking, we fix diagnostic OD cases via a ground-truth corridor audit (Sec.~\ref{sec:experiments}).
We observe two characteristic symptoms of topological multi-modality: (i) deterministic coordinate regression collapses to overly direct ``average'' routes, and (ii) stochastic generation can avoid hard collapse yet still blurs corridor-level alternatives into implausible mixtures.
Together, these failures suggest that the bottleneck is not simply model capacity, but \emph{mode interference} when corridor-level decisions are represented only implicitly in a continuous coordinate output space.
Critically, our ablations show that corridor diversity must be (i) created as an explicit \emph{topological commitment} at the decision stage and (ii) \emph{preserved} during execution: a strong continuous refiner can overwrite corridor commitments, while a commitment-preserving refinement maintains coverage.

\begin{figure}[t]
  \centering
  \includegraphics[width=0.95\linewidth]{figures/fig_mode_collapse.pdf}
  \caption{Diagnostic visualization on a fixed multi-modal OD case. Deterministic L2 regression collapses to an averaged route; end-to-end diffusion avoids hard collapse yet exhibits corridor-level interference; a diffusion-based waypoint decision model recovers both corridor alternatives by explicit multi-modal sampling, and a cascaded generator can preserve this commitment when execution is applied as a \emph{weak refinement}.}
  \label{fig:mode-collapse}
\end{figure}

\paragraph{CascadeTraj: bridging semantic intent and physical execution.}
We propose \textbf{CascadeTraj}, a cascaded generative framework for route generation.
At the \emph{decision} stage, the model samples a sparse route skeleton (e.g., waypoints) that represents discrete corridor-level alternatives and preserves long-horizon intent consistency.
At the \emph{execution} stage, the model generates a full-resolution continuous route conditioned on the skeleton, capturing local geometry and motion patterns.
This decomposition turns a brittle one-shot generation problem into two conditional problems with clearer responsibilities: (i) committing to a coherent topological plan, and (ii) executing a physically plausible realization.

\paragraph{Soft priors and optional census guidance.}
To encourage feasibility without collapsing support, we incorporate map information as \emph{soft} priors (e.g., road-proximity fields) rather than hard output masking.
Our decision stage also admits optional spatial semantics (e.g., POI/land-use covariates, or census-derived attributes) as conditioning signals.
These covariates can steer route preferences and, in our ablations, can further improve corridor-level diversity, but they are \emph{not required} for the core mechanism in this paper.
While macro-level statistics can provide useful priors for controllability, this paper focuses on the generative architecture for the \emph{route itself}; census integration is explored only as an ablation to demonstrate extensibility.

\paragraph{Contributions.}
\begin{itemize}
  \item \textbf{Finding 1 (Diagnosis):} we provide evidence that topological multi-modality induces mode interference and destructive averaging in end-to-end long-horizon route generators, which cannot be fixed by naively scaling continuous coordinate models.
  \item \textbf{Finding 2 (Mechanism):} we find that corridor-level diversity must be represented explicitly as a topological commitment and then preserved through composition: a waypoint bottleneck only helps when the skeleton is sampled multi-modally, and execution must act as a commitment-preserving refinement rather than overwriting corridor choices.
  \item \textbf{Robustness:} we incorporate map information as \emph{soft} priors (instead of hard masking) to encourage feasibility without truncating support, and treat additional macro semantics (e.g., census covariates) strictly as optional ablations for extensibility.
\end{itemize}
